\subsection{The projection space is load-bearing: its size, its nonlinearity, and its conditioning}
\label{sec:projection-space}

Everything above treats the rectangular projection $P_m$ as fixed plumbing: a
$d_{\text{shared}}\times D_m$ matrix carrying each base into the shared routing space. That
assumption turned out to be the largest untested surface in the study, and testing it changes what
the headline numbers mean. We report it as a separate line of work because it does not merely
refine an existing claim---it relocates the mechanism.

\paragraph{The projection in every fusion experiment was frozen and random.}
In the scale-heterogeneous configuration of \S\ref{sec:scale-hetero-exp} and in the
architectural configuration of \S\ref{sec:fusion-exp}, every $P_m$ is orthogonally initialized and
then \emph{frozen}, the per-expert feed-forward blocks are frozen as well, and the only trainable
tensors are the shared router and three per-expert affine pairs ($\approx 2.3$K parameters in
total). The margins reported for both lines were therefore obtained inside a shared space whose
geometry was random. Any claim of the form ``fusion beats its parts'' must be read against that
fact.

\paragraph{How large does the space need to be?}
A training-free diagnostic (no gradients) stacks the bases' token representations, takes the
covariance spectrum, and measures cross-base redundancy with linear centered kernel alignment
(CKA) against a row-permutation null. On a controlled corpus of $8$ semantic classes built from
slot combinations---so that the training and held-out splits share no sentence at all---the union of
the three bases occupies $4$--$19$ dimensions at the $95$--$99\%$ variance level, against the
$d_{\text{shared}}{=}256$ actually used; the pairwise CKA is far above its permutation null. Two
things follow. The space is provisioned for roughly an order of magnitude more directions than the
data occupy, and the redundancy that would justify a small space is real---but it is invisible to
a linear method: the joint principal subspace of the union, which is the best a linear
construction can do, leaves the three bases in nearly disjoint directions (footprint overlap
$<0.06$).

\paragraph{Nonlinearity, not capacity, closes the gap.}
We then run a $2\times2$ factorial over the two axes that the diagnosis implicates---nonlinear
versus linear lifting, and static versus content-conditioned lifting---crossed with a sample-size
axis ($32$, $256$, $1024$ examples per base), with the hidden width of the nonlinear lift solved
so that the two arms differ in fewer than $0.3\%$ of parameters, and with a fixed minibatch and
step count so that the sample-size axis measures data diversity rather than compute. The
nonlinear arm reduces held-out fusion error by $4.7$--$9.9\times$ relative to the
parameter-matched linear arm, and does so with roughly half the overfitting gap; the linear arm's
advantage instead comes from sample size, and shrinks as the nonlinear arm saturates. A
random-projection construction is $38\times$ worse than the trivial zero predictor in absolute
error. The redundancy the diagnostic located is therefore recoverable, but only through a nonlinear
map.

\paragraph{Conditioning needs a nonlinear base, and needs data.}
Conditioning the projection on a content-derived vector (a hypernetwork emitting a rank-16
modulation) is useless on a linear base---the conditioned and static linear arms agree to within
noise at all three sample sizes---and becomes useful on top of the nonlinear lift
($@NL_COND_S1024@$ versus $@NL_S1024@$ at $1024$ examples). Permuting the conditions across samples
degrades held-out error by $+24.6\%$, $+47.0\%$ and $+50.3\%$ as the sample size grows, so the
conditioning is used rather than decorative. It also requires data: at $32$ examples the
conditioned nonlinear arm is \emph{worse} than the unconditioned one.

\paragraph{Auditing the margin itself.}
The finding above makes the headline margin suspect, and we test it directly rather than by
argument. If the frozen random space inflates the fused-over-single-expert margin, then in that
configuration the fused output should carry the \emph{identity of its source base} rather than the
semantics of the input. We test this with two linear probes on the held-out fused output: one
predicts which of the three bases produced the sample, the other predicts the target class. In the
frozen random configuration the margin over the best single expert is $@A_GAIN@$ (held-out fusion
error $@A_HELD@$, against a trivial zero-predictor error of $0.125$), and the probes return
$@A_SRC@$ for source identity and only $@A_CLS@$ for the target class, while the projected inputs
themselves identify their source with probability $@A_IDENT@$. The fused output is, in effect, a
source detector: the router's task is trivial in that space, and the margin over a single expert
measures source detection and dispatch rather than semantic fusion. Across the audited arms the
margin over the best single expert is $@A_GAIN@$ (frozen random), $@B_GAIN@$ (nonlinear, frozen
random), $@C_GAIN@$ (nonlinear, trained), $@D_GAIN@$ (nonlinear, proxy-aligned, frozen) and
$@E_GAIN@$ (nonlinear, proxy-aligned, trained), with held-out errors of $@A_HELD@$, $@C_HELD@$ and
$@E_HELD@$ respectively.

\paragraph{What we therefore claim.}
We cannot attribute the earlier margins to fusion. They are attributable to a shared space that
was never optimized, and they do not persist once the space is learned: with a meaningful
(nonlinear, trained) space, a single expert with the same space solves the task about as well as
the fused model does. The defensible claim is the weaker and more useful one: \emph{a shared space
must be meaningful for sublayer fusion to work at all}---in a random space the fused output is
$38\times$ worse than predicting zeros---and once it is meaningful, fusing heterogeneous bases
adds little over using one of them well. We report the negative result with the same prominence
as the positive ones, and we mark the audit as pending a fully independent re-derivation.
